% SPDX-License-Identifier: CC-BY-SA-4.0
% Author: Matthieu Perrin

\begingroup

\SetKwBroadcast{STA}{sta}
\SetKwBroadcast{RB}{rb}
\SetKwBroadcast{URB}{urb}
\SetKwBroadcast{FIFO}{fifo}
\SetKwMessage{MSG}{msg}
\SetKwData{Clock}{clock}
\SetKwData{Delivered}{delivered}
\SetKwData{Tick}{tick}
\SetKwData{Now}{now}
\SetKwData{Update}{update}
\SetKwVariable{Payload}{payload}
\SetKwVariable{Date}{date}
\SetKwVariable{SN}{sn}

\begin{frame}{Implémentation de FIFO Broadcast}

  \ob<-2>[top]{
    \begin{algorithm}[H]
      \LVariables{\At $p_i$}{
        $\Delivered_i \leftarrow [0, ..., 0]$ \tcp*[f]{Compteurs des messages délivrés}
      }
      \Method{$\FIFO.\Broadcast(\Payload)$ \At $p_i$}{
        \nl \STA.\SBroadcast $\alert{\MSG(\Payload, \Delivered_i[i])}$\\
      }
      
      \Upon{\STA.\Deliver $\alert{\MSG(\Payload, \SN)}$ \From $p_j$ \At $p_i$}{
        \nl \Wait $\Delivered_j[i] = \SN$\\
        \nl $\Delivered_j[i] \leftarrow \SN+1$\\
        \nl \Trigger \FIFO.\Deliver $\Payload$ \From $p_i$
      }
    \end{algorithm}
  }

  \on<3>[top]{
    \begin{algorithm}[H]
      \LVariables{\At $p_i$}{
        \Example{$\Clock_i[0, ..., 0]$} \tcp*[f]{$\Delivered_i[1..n]$}
      }
      \Method{$\FIFO.\Broadcast(\Payload)$ \At $p_i$}{
        \nl\STA.\SBroadcast $\alert{\MSG(\Payload, \Example{\Clock_i[i].\Now()})}$ \tcp*[r]{$\Delivered_i[i]$}
      }
      \Upon{\STA.\Deliver $\alert{\MSG(\Payload, \Date)}$ \From $p_i$ \At $p_j$}{
        \nl \Wait $\Example{\Clock_j[i].\Now() \ge \Date}$ \tcp*[r]{$\Delivered_j[i]$}
        \nl \Example{$\Clock_j[i].\Update()$} \tcp*[r]{$\Delivered_j[i] \leftarrow \SN+1$}
        \nl \Trigger \FIFO.\Deliver $\Payload$ \From $p_i$
      }
    \end{algorithm}
  }

  \obExampleBlock<1>[anchor=north, y=-2mm]{Exemple d'exécution}{}
  \ob<1>[y=-27mm]{
    \begin{tikzpicture}[y=12mm]
      \draw[process] (0,1) node[left]{$p_1$}                                            to (10,1);
      \draw[process] (0,0) node[left]{$p_2$} node[replica below]{$\Delivered_2[1] = 0$} to (10,0);

      \draw[alert]   (1,1) node (m1) {} node[above] {$m$} ;
      \draw[alert, message]     (m1.center) to[bend left]   (2,1) ;
      \draw[alert, message]     (m1.center) to node[swap]{$sn=0$} (4.5,0) node[replica below]{$1$};

      \draw[example] (3,1) node (m2) {} node[above] {$m'$} ;
      \draw[example, message]   (m2.center) to[bend left]   (4,1) ;
      \draw[example, message]   (m2.center) to node{$sn=1$} (3.5,.05) to[bend left] (5.5,0) node[replica below]{$2$} ;
    \end{tikzpicture}
  }

  \onBlock<2->[anchor=north, y=-2mm]{Remarques}{
    \begin{itemize}
    \item FIFO-Ordering : pour tout $j$, $\Delivered_i[j]$ = prochain $sn$ attendu de $p_j$.
    \item On peut renforcer les propriétés en remplaçant $\STA$ par $\RB$ ou $\URB$
    \item En pratique, \Wait est implémenté par des buffers
    \item<3> Pour tout processus $p_j$, \example{$\Delivered_i[j]$} implémente une horloge logique
      \begin{itemize}
      \item Le temps logique est l'ordre de processus
      \item Les dates sont des entiers (horloge scalaire)
      \end{itemize}
    \end{itemize}
  }
  
\end{frame}

\endgroup
\endinput
